\Article{23}
The work by the distinguished \textsc{de Krayenhof}, \textit{Précis historique des opérations trigonométriques faites en Hollande}, will furnish our first example. We shall subject to adjustment that part of the network of triangles which lies among the nine points Harlingen, Sneek, Oldeholtpade, Ballum, Leeuwarden, Dockum, Drachten, Oosterwolde, and Gröningen. Between these points nine triangles are formed, designated in that work by the numbers 121, 122, 123, 124, 125, 127, 128, 131, and 132. Their angles, distinguished by the indices specified by us, were observed as follows, according to the table on pp. 77–81:
\SourceLeaf{088}
\begin{center}
\begin{tabular}{r l l r r r}
\multicolumn{6}{c}{Triangle 121.}\\[5pt]
0.&Harlingen&. . . . .&\(50^0\)&\(58'\)&\(15''238\)\\
1.&Leeuwarden&. . .&82&47&15,351\\
2.&Ballum&. . . . . . .&46&14&27,202
\end{tabular}
\end{center}
\begin{center}
\begin{tabular}{r l l r r r}
\multicolumn{6}{c}{Triangle 122.}\\[5pt]
3.&Harlingen&. . . . .&51&5&39,717\\
4.&Sneek&. . . . . . . .&70&48&33,445\\
5.&Leeuwarden&. . .&58&5&48,707
\end{tabular}
\end{center}
\begin{center}
\begin{tabular}{r l l r r r}
\multicolumn{6}{c}{Triangle 123.}\\[5pt]
6.&Sneek&. . . . . . . .&49&30&40,051\\
7.&Drachten&. . . . . .&42&52&59,382\\
8.&Leeuwarden&. . .&87&36&21,057
\end{tabular}
\end{center}
\begin{center}
\begin{tabular}{r l l r r r}
\multicolumn{6}{c}{Triangle 124.}\\[5pt]
9.&Sneek&. . . . . . . .&45&36&7,492\\
10.&Oldeholtpade&. . .&67&52&0,048\\
11.&Drachten&. . . . .&66&31&56,513
\end{tabular}
\end{center}
\begin{center}
\begin{tabular}{r l l r r r}
\multicolumn{6}{c}{Triangle 125.}\\[5pt]
12.&Drachten&. . . . . .&53&55&24,745\\
13.&Oldeholtpade&. . .&47&48&52,580\\
14.&Oosterwolde&. . .&78&15&42,347
\end{tabular}
\end{center}
\begin{center}
\begin{tabular}{r l l r r r}
\multicolumn{6}{c}{Triangle 127.}\\[5pt]
15.&Leeuwarden&. . .&59&24&0,645\\
16.&Dockum&. . . . . .&76&34&9,021\\
17.&Ballum&. . . . . . .&44&1&51,040
\end{tabular}
\end{center}
\begin{center}
\begin{tabular}{r l l r r r}
\multicolumn{6}{c}{Triangle 128.}\\[5pt]
18.&Leeuwarden&. . .&72&6&32,043\\
19.&Drachten&. . . . . .&46&53&27,163\\
20.&Dockum&. . . . . .&61&0&4,494
\end{tabular}
\end{center}
\begin{center}
\begin{tabular}{r l l r r r}
\multicolumn{6}{c}{Triangle 131.}\\[5pt]
21.&Dockum&. . . . . .&57&1&55,292\\
22.&Drachten&. . . . . .&83&33&14,515\\
23.&Gröningen&. . . . .&39&24&52,397
\end{tabular}
\end{center}
\SourceLeaf{089}
\begin{center}
\begin{tabular}{r l l r r r}
\multicolumn{6}{c}{Triangle 132.}\\[5pt]
24.&Oosterwolde&. . .&\(81^0\)&\(54'\)&\(17''447\)\\
25.&Gröningen&. . . . .&31&52&46,094\\
26.&Drachten&. . . . .&66&12&57.246
\end{tabular}
\end{center}

An examination of the connections among these triangles shows that among the 27 angles whose approximate values have become known through observation there are 13 conditional equations: namely, two of the first kind, nine of the second, and two of the third. It is unnecessary to write all these equations here in their finite form, since the calculations require only the quantities denoted in the general theory by \(\mathfrak A,a,a',a''\), etc., \(\mathfrak B,b,b',b''\), etc. We therefore write at once, in their place, the equations designated above by (13), which exhibit those quantities; instead of the symbols \(\varepsilon,\varepsilon',\varepsilon''\), etc., we shall simply write \((0),(1).\ (2)\), etc.

In this way, the following equations correspond to the two conditional equations of the first kind:
\[
\begin{aligned}
(1)+(5)+(8)+(15)+(18)&=-2''197\\
(7)+(11)+(12)+(19)+(22)+(26)&=-0''436
\end{aligned}
\]

We next find the spheroidal excesses of the nine triangles to be: \(1''749;\allowbreak1''147;\allowbreak1''243;\allowbreak1''698;\allowbreak0''873;\allowbreak1''167;\allowbreak1''104;\allowbreak2''161;\allowbreak1''403\). Thus the first conditional equation of the second kind is the following\footnote{. In this example we prefer to express the indices with Arabic numerals.}: \(v^{(0)}+v^{(1)}+v^{(2)}-180^0\,0'\,1''749=0\), and the others are obtained in the same way; hence we have the following nine equations:
\[
\begin{aligned}
(0)+(1)+(2)&=-3''958\\
(3)+(4)+(5)&=+0{,}722\\
(6)+(7)+(8)&=-0{,}753\\
(9)+(10)+(11)&=+2{,}355\\
(12)+(13)+(14)&=-1{,}201\\
(15)+(16)+(17)&=-0{,}461\\
(18)+(19)+(20)&=+2{,}596\\
(21)+(22)+(23)&=+0{,}043\\
(24)+(25)+(26)&=-0{,}616
\end{aligned}
\]
The conditional equations of the third kind are more conveniently presented in logarithmic form; the first is therefore
\SourceLeaf{090}
\[
\begin{aligned}
&\log\sin(v^{(0)}-0''583)-\log\sin(v^{(2)}-0''583)\\
&\quad-\log\sin(v^{(3)}-0''382)+\log\sin(v^{(4)}-0''382)\\
&\quad-\log\sin(v^{(6)}-0''414)+\log\sin(v^{(7)}-0''414)\\
&\quad-\log\sin(v^{(16)}-0''389)+\log\sin(v^{(17)}-0''389)\\
&\quad-\log\sin(v^{(19)}-0''368)+\log\sin(v^{(20)}-0''368)=0
\end{aligned}
\]
It seems unnecessary to write the other one in finite form. The following equations correspond to these two equations; each coefficient refers to the seventh figure of the Briggsian logarithms:
\[
\begin{aligned}
&17{,}068(0)-20{,}174(2)-16{,}993(3)+7{,}328(4)\\
&\quad-17{,}976(6)+22{,}672(7)-5{,}028(16)+21{,}780(17)\\
&\quad-19{,}710(19)+11{,}671(20)=-371\\[5pt]
&17{,}976(6)-0{,}880(8)-20{,}617(9)+8{,}564(10)\\
&\quad-19{,}082(13)+4{,}375(14)+6{,}798(18)-11{,}671(20)\\
&\quad+13{,}657(21)-25{,}620(23)-2{,}995(24)+33{,}854(25)=+370
\end{aligned}
\]

Since no reason has been indicated for assigning unequal weights to the observations, we shall set \(p^{(0)}=p^{(1)}=p^{(2)}\text{etc.}=1\). Denoting the correlates of the conditional equations, in the order in which we have presented the corresponding equations, by \(A,\allowbreak B,\allowbreak C,\allowbreak D,\allowbreak E,\allowbreak F,\allowbreak G,\allowbreak H,\allowbreak I,\allowbreak K,\allowbreak L,\allowbreak M,\allowbreak N\), the following equations result for determining them:
\[
\begin{aligned}
-2''197&=5A+C+D+E+H+I+5{,}917N\\
-0{,}436&=6B+E+F+G+I+K+L+2{,}962M\\
-3{,}958&=A+3C-3{,}106M\\
+0{,}722&=A+3D-9{,}665M\\
-0{,}753&=A+B+3E+4{,}696M+17{,}096N\\
+2{,}355&=B+3F-12{,}053N\\
-1{,}201&=B+3G-14{,}707N\\
-0{,}461&=A+3H+16{,}752M\\
+2{,}596&=A+B+3I-8{,}039M-4{,}874N\\
+0{,}043&=B+3K-11{,}963N\\
-0{,}616&=B+3L+30{,}859N\\
-371&=+2{,}962B-3{,}106C-9{,}665D+4{,}696E\\
&\quad+16{,}752H-8{,}039I+2902{,}27M-459{,}33N\\
+370&=+5{,}917A+17{,}096E-12{,}053F-14{,}707G\\
&\quad-4{,}874I-11{,}963K+30{,}859L-459{,}33M+3385{,}96N
\end{aligned}
\]
\SourceLeaf{091}

By elimination we obtain:
\[
\begin{array}{r@{\;}c@{\;}l|r@{\;}c@{\;}l}
A&=&-0{,}598&H&=&+0{,}659\\
B&=&-0{,}255&I&=&+1{,}050\\
C&=&-1{,}234&K&=&+0{,}577\\
D&=&+0{,}086&L&=&-1{,}351\\
E&=&-0{,}477&M&=&-0{,}109792\\
F&=&+1{,}351&N&=&+0{,}119681\\
G&=&+0{,}271&&&
\end{array}
\]
Finally, the most probable errors are given by the formulas
\[
\begin{aligned}
(0)&=C+17{,}068M\\
(1)&=A+C\\
(2)&=C-20{,}174M\\
(3)&=D-16{,}993M
\end{aligned}
\]
etc., whence we obtain the following numerical values; for comparison, we append, with their signs changed, the corrections applied by the distinguished \textsc{de Krayenhof} to the observations:
\[
\begin{array}{r@{\;}c@{\;}r|r||r@{\;}c@{\;}r|r}
&&&\text{\textsc{de Kr.}}&&&&\text{\textsc{de Kr.}}\\[4pt]
(0)&=&-3''108&-2''090&(14)&=&+0''795&+2''400\\
(1)&=&-1{,}832&+0{,}116&(15)&=&+0{,}061&+1{,}273\\
(2)&=&+0{,}981&-1{,}982&(16)&=&+1{,}211&+5{,}945\\
(3)&=&+1{,}952&+1{,}722&(17)&=&-1{,}732&-7{,}674\\
(4)&=&-0.719&+2{,}848&(18)&=&+1{,}265&+1{,}876\\
(5)&=&-0{,}512&-3{,}848&(19)&=&+2{,}959&+6{,}251\\
(6)&=&+3{,}648&-0{,}137&(20)&=&-1{,}628&-5{,}530\\
(7)&=&-3{,}221&+1{,}000&(21)&=&+2{,}211&+3{,}486\\
(8)&=&-1{,}180&-1{,}614&(22)&=&+0{,}322&-3{,}454\\
(9)&=&-1{,}116&0&(23)&=&-2{,}489&0\\
(10)&=&+2{,}376&+5.928&(24)&=&-1{,}709&+0{,}400\\
(11)&=&+1{,}096&-3{,}570&(25)&=&+2{,}701&+2.054\\
(12)&=&+0{,}016&+2{,}414&(26)&=&-1{,}606&-3{,}077\\
(13)&=&-2{,}013&-6{,}014&&&&
\end{array}
\]
\SourceLeaf{092}

{\rightskip=0pt plus 1em The sum of the squares of our adjustments is found to be \(=97{,}8845\). Hence the mean error, insofar as it can be inferred from the 27 observed angles, is\par}
\[
=\BareRoot{\frac{97{,}8845}{13}}=2''7440
\]

The sum of the squares of the changes that the distinguished \textsc{de Krayenhof} himself applied to the observed angles is found to be \(=341{,}4201\).

